% chapters/02_linux_docker.tex
\chapter{Linux and Docker: the foundation of the stack}
\label{ch:linux_docker}

\section{Goal (what you will achieve)}
This chapter builds the foundation for everything else.

By the end, you will have:
\begin{itemize}[leftmargin=*]
  \item a Linux environment ready to run services continuously,
  \item Docker installed and working,
  \item Docker Compose available for multi-container applications,
  \item a clean folder layout for persistent data (so upgrades do not destroy your files),
  \item a simple way to verify your system is healthy before installing the applications.
\end{itemize}

\section{Before we start: two valid ways to get Linux}
When beginners hear ``install Docker on Linux'', a common confusion is:
\emph{``But I am on Windows---where is this Linux terminal you keep mentioning?''}

You need \textbf{some Linux environment} before you can run the Linux commands in this manual.
You can get that in two common ways:

\subsection{Path A (recommended): install Linux directly on the machine}
This is the most reliable approach for a real server:
\begin{itemize}[leftmargin=*]
  \item your machine boots into Linux,
  \item you open the Linux terminal directly,
  \item services run in the background like a proper server.
\end{itemize}

\subsection{Path B (beginner-friendly on Windows): use WSL2}
WSL2 (Windows Subsystem for Linux) lets you run a real Linux distribution inside Windows.
You will open a Linux terminal \emph{from Windows}, but the commands still run \emph{inside Linux}.

This is a good way to learn and to prototype, but for a long-term always-on server,
a full Linux install is usually simpler.

\section{Important concept: what is a ``terminal'' and what is \texttt{sudo}?}
A \textbf{terminal} is a text-based interface to the operating system. You type commands, the system executes them.

\texttt{sudo} means ``run as administrator'' \emph{inside Linux}.
\begin{itemize}[leftmargin=*]
  \item On Linux, many system changes (installing software, editing protected files) require administrator privileges.
  \item \texttt{sudo} temporarily grants those privileges when you trust the command.
  \item On Windows, \texttt{sudo} is not a standard command (unless you are in a Linux environment such as WSL).
\end{itemize}

So if you ever wonder ``why does sudo work on my Windows PC?'', the answer is:
\textbf{because you are actually inside Linux (WSL), not executing Windows commands.}

\section{Step 1 --- Get a Linux terminal (choose A or B)}
\subsection{Step 1A: Linux installed directly (Ubuntu Server / Debian)}
At the end of the Linux installation you can:
\begin{itemize}[leftmargin=*]
  \item log in on the machine itself (keyboard + screen), or
  \item log in remotely using SSH (recommended).
\end{itemize}

If you are on the machine: you will see a login prompt. After login, you are already in a terminal.

If you are connecting remotely (SSH): from another computer you use an SSH client.
On Windows, you can open \textbf{PowerShell} and type:
\begin{lstlisting}
ssh username@SERVER_IP
\end{lstlisting}
After logging in, you are inside the Linux terminal of the server.

\subsection{Step 1B: Windows + WSL2 (Ubuntu inside Windows)}
If you want the beginner path on Windows, do this:
\begin{itemize}[leftmargin=*]
  \item Install WSL2 and an Ubuntu distribution.
  \item Open the Ubuntu app (this is your Linux terminal).
  \item Run all Linux commands from that Ubuntu window.
\end{itemize}

In practice, once WSL2 is installed, you open \textbf{Ubuntu} from the Start menu.
That window is the Linux terminal.

\paragraph{Note.}
This manual does not assume you already know WSL2. If you choose this path,
you can treat WSL2 as ``a Linux box living inside Windows''.

\section{Step 1 (continued) --- Update Linux}
Once you have \emph{a Linux terminal open} (either on a Linux machine, via SSH, or via WSL),
run the following commands.

\subsection{1.1 Update package lists}
\begin{lstlisting}
sudo apt update
\end{lstlisting}
\textbf{What this does:} downloads the latest list of available software packages from your configured repositories. \\
\textbf{Why you do it:} ensures you install up-to-date packages and security fixes.

\subsection{1.2 Upgrade installed packages}
\begin{lstlisting}
sudo apt upgrade -y
\end{lstlisting}
\textbf{What this does:} upgrades installed packages to newer versions. \\
\textbf{Why you do it:} reduces bugs and security issues before you install Docker. \\
\textbf{About \texttt{-y}:} automatically answers ``yes'' to prompts (useful in step-by-step guides).

\subsection{1.3 Install basic tools we will use later (recommended)}
\begin{lstlisting}
sudo apt install -y curl ca-certificates gnupg lsb-release unzip
\end{lstlisting}
\textbf{What these are:}
\begin{itemize}[leftmargin=*]
  \item \texttt{curl}: downloads data from the internet (we use it to fetch Docker keys and scripts),
  \item \texttt{ca-certificates}: ensures HTTPS downloads are verified correctly,
  \item \texttt{gnupg}: used to verify signed packages (security),
  \item \texttt{lsb-release}: helps detect your Linux version codename automatically,
  \item \texttt{unzip}: needed sometimes for downloaded archives.
\end{itemize}

\subsection{Quick sanity check: do you have \texttt{sudo} access?}
Try:
\begin{lstlisting}
sudo whoami
\end{lstlisting}
If it prints \texttt{root}, your user can run admin commands (good).
If it fails, your user is not allowed to use \texttt{sudo} and you must fix that first
(typically by using the admin account created during Linux installation).

\section{Step 2 --- Create a clean data layout}
A common beginner mistake is to store important data \emph{inside} containers.
Do \textbf{not} do that.

\subsection*{The idea (in one sentence)}
Containers are meant to be disposable; your data should live in \emph{normal folders on disk} that you control.

\subsection*{Choose a single root folder for persistent data}
Instead of scattering files across the system, choose one directory that will hold:
\begin{itemize}[leftmargin=*]
  \item your application configuration (Compose files),
  \item your databases and indexes,
  \item your photo/music libraries,
  \item backups.
\end{itemize}

Two common options are:
\begin{itemize}[leftmargin=*]
  \item \texttt{/srv} (very common on servers),
  \item \texttt{/opt} (also used for application data).
\end{itemize}

In this manual we will use \texttt{/srv}.

\subsection*{Create the folders}
\begin{lstlisting}
sudo mkdir -p /srv/{docker,photos,music,drive,ai,backups}
sudo mkdir -p /srv/docker/compose
\end{lstlisting}

\textbf{What this does:}
\begin{itemize}[leftmargin=*]
  \item \texttt{mkdir} creates directories,
  \item \texttt{-p} means ``create parent folders if needed and do not complain if the folder already exists'',
  \item the \texttt{\{...\}} syntax is a shortcut to create multiple folders in one command.
\end{itemize}

\textbf{Why you do it:} you get a predictable, documented structure that makes your stack easier to maintain and back up.

\subsection*{Give your user ownership}
\begin{lstlisting}
sudo chown -R $USER:$USER /srv
\end{lstlisting}

\textbf{What this does:}
\begin{itemize}[leftmargin=*]
  \item \texttt{chown} changes ownership of files/folders,
  \item \texttt{-R} means ``recursive'' (apply to everything inside),
  \item \texttt{\$USER} expands to your current username.
\end{itemize}

\textbf{Why you do it:} so you can edit configuration and manage volumes without constantly using \texttt{sudo}.

\subsection{Why this matters}
When you upgrade or recreate containers, your data remains safe in \texttt{/srv}.
Containers can be disposable; your files should not be.

\section{Step 3 --- Install Docker Engine}
The safest approach is to install Docker from Docker's official repository.
This matters because:
\begin{itemize}[leftmargin=*]
  \item you get security updates and bug fixes,
  \item you avoid outdated versions from generic repositories,
  \item you can reproduce the same setup on another machine.
\end{itemize}

\subsection{3.1 Add Docker's official GPG key}
\begin{lstlisting}
sudo install -m 0755 -d /etc/apt/keyrings
curl -fsSL https://download.docker.com/linux/ubuntu/gpg | sudo gpg --dearmor -o /etc/apt/keyrings/docker.gpg
sudo chmod a+r /etc/apt/keyrings/docker.gpg
\end{lstlisting}

\textbf{What this does (line by line):}
\begin{itemize}[leftmargin=*]
  \item \texttt{install -d ...} creates a system folder for trusted repository keys,
  \item \texttt{curl ... | gpg --dearmor ...} downloads Docker's signing key and stores it in a format APT can use,
  \item \texttt{chmod a+r ...} makes the key readable by the package manager.
\end{itemize}

\textbf{Why you do it:} APT uses this key to verify that Docker packages are authentic (not tampered with).

\subsection{3.2 Add the Docker repository}
\textbf{Important: choose the correct block for your distribution.}

\subsubsection*{Ubuntu}
\begin{lstlisting}
echo \
  "deb [arch=$(dpkg --print-architecture) signed-by=/etc/apt/keyrings/docker.gpg] https://download.docker.com/linux/ubuntu \
  $(. /etc/os-release && echo "$VERSION_CODENAME") stable" | \
  sudo tee /etc/apt/sources.list.d/docker.list > /dev/null
\end{lstlisting}

\subsubsection*{Debian}
\begin{lstlisting}
echo \
  "deb [arch=$(dpkg --print-architecture) signed-by=/etc/apt/keyrings/docker.gpg] https://download.docker.com/linux/debian \
  $(. /etc/os-release && echo "$VERSION_CODENAME") stable" | \
  sudo tee /etc/apt/sources.list.d/docker.list > /dev/null
\end{lstlisting}

\textbf{What this does:}
\begin{itemize}[leftmargin=*]
  \item writes a new ``software source'' file for APT,
  \item automatically detects your CPU architecture (e.g., amd64),
  \item automatically detects your distro codename (e.g., \texttt{jammy}, \texttt{bookworm}),
  \item tells APT to trust packages signed by the Docker key you installed.
\end{itemize}

\textbf{Why you do it:} so \texttt{apt} can download Docker directly from the official source.

\subsection{3.3 Install Docker Engine and plugins}
\begin{lstlisting}
sudo apt update
sudo apt install -y docker-ce docker-ce-cli containerd.io docker-buildx-plugin docker-compose-plugin
\end{lstlisting}

\textbf{What this does:}
\begin{itemize}[leftmargin=*]
  \item \texttt{apt update} refreshes package lists (now including Docker's repo),
  \item installs Docker Engine (\texttt{docker-ce}) and its command-line interface (\texttt{docker-ce-cli}),
  \item installs \texttt{containerd} (the container runtime Docker uses under the hood),
  \item installs Buildx (modern build features),
  \item installs the Compose plugin (so you can use \texttt{docker compose ...}).
\end{itemize}

\textbf{Why you do it:} this gives you a modern Docker + Compose setup that matches current documentation.

\subsection{3.4 Enable and start Docker}
\begin{lstlisting}
sudo systemctl enable docker
sudo systemctl start docker
\end{lstlisting}

\textbf{What this does:}
\begin{itemize}[leftmargin=*]
  \item \texttt{enable} makes Docker start automatically at boot,
  \item \texttt{start} starts Docker right now.
\end{itemize}

\textbf{Why you do it:} your services should survive reboots.

\paragraph{WSL note.}
If you are on WSL2, \texttt{systemctl} may not work depending on your setup.
In that case, Docker is typically used via Docker Desktop on Windows, or via a WSL configuration that supports systemd.

\section{Step 4 --- Allow running Docker without sudo (recommended)}
By default, Docker requires \texttt{sudo}. For day-to-day use, add your user to the \texttt{docker} group:

\begin{lstlisting}
sudo usermod -aG docker $USER
\end{lstlisting}

\textbf{What this does:}
\begin{itemize}[leftmargin=*]
  \item adds your user to the \texttt{docker} group,
  \item users in that group can talk to the Docker daemon without \texttt{sudo}.
\end{itemize}

\textbf{Why you do it:} it makes everyday commands less annoying.

Then log out and log back in (or reboot) so the group membership applies.

\section{Step 5 --- Verify the installation}
\subsection{Docker version}
\begin{lstlisting}
docker --version
\end{lstlisting}
\textbf{What this does:} prints the installed Docker version. \\
\textbf{Why you do it:} confirms the command is available.

\subsection{Compose version}
\begin{lstlisting}
docker compose version
\end{lstlisting}
\textbf{What this does:} prints the Compose plugin version. \\
\textbf{Why you do it:} confirms you can run multi-service stacks.

\subsection{Hello-world test}
\begin{lstlisting}
docker run --rm hello-world
\end{lstlisting}
\textbf{What this does:}
\begin{itemize}[leftmargin=*]
  \item downloads a tiny test image (if needed),
  \item runs it in a container,
  \item \texttt{--rm} deletes the container after it exits (keeps things clean).
\end{itemize}

\textbf{Why you do it:} it proves the full chain works: download images + run containers.

\section{Step 6 --- (Optional) Basic firewall and SSH hygiene}
If you are only using the services inside your home network or behind a VPN, you can keep things simple.
If your machine is exposed to the internet, do \textbf{not} skip this.

\subsection{Enable UFW (simple firewall)}
\begin{lstlisting}
sudo apt install -y ufw
sudo ufw default deny incoming
sudo ufw default allow outgoing
sudo ufw allow OpenSSH
sudo ufw enable
sudo ufw status
\end{lstlisting}

\textbf{What this does:}
\begin{itemize}[leftmargin=*]
  \item installs UFW (Uncomplicated Firewall),
  \item blocks all incoming traffic by default (safer),
  \item allows outgoing traffic (normal),
  \item allows SSH so you do not lock yourself out,
  \item enables the firewall and shows its status.
\end{itemize}

\textbf{Why you do it:} it reduces the attack surface of your server.

\subsection{SSH: disable password login (recommended)}
Once you can log in via SSH keys, disable password login by editing the SSH configuration file.

\begin{lstlisting}
sudo nano /etc/ssh/sshd_config
\end{lstlisting}

\textbf{What this does:} opens the SSH server configuration in a text editor. \\
\textbf{Why you do it:} to prevent password brute-force attacks.

Find the line (or add it) and set:

\begin{lstlisting}
PasswordAuthentication no
\end{lstlisting}

Then restart SSH:

\begin{lstlisting}
sudo systemctl restart ssh
\end{lstlisting}

\textbf{What this does:} restarts the SSH service so changes take effect. \\
\textbf{Why you do it:} config changes do nothing until the service reloads/restarts.

\section{Verification checklist}
Before moving on, confirm:
\begin{itemize}[leftmargin=*]
  \item \texttt{docker run --rm hello-world} works,
  \item \texttt{docker compose version} prints a version,
  \item you have a persistent folder structure under \texttt{/srv},
  \item you understand the core idea: containers are disposable; data is not.
\end{itemize}

\section{Troubleshooting (common issues)}
\subsection{``Permission denied'' when running Docker}
If you see permission errors without \texttt{sudo}, either:
\begin{itemize}[leftmargin=*]
  \item you did not add your user to the \texttt{docker} group, or
  \item you did, but you did not log out/in (group not applied yet).
\end{itemize}

Fix:
\begin{lstlisting}
sudo usermod -aG docker $USER
\end{lstlisting}
Then log out and log back in.

\subsection{Docker service not running}
Check status:
\begin{lstlisting}
sudo systemctl status docker
\end{lstlisting}
Start it:
\begin{lstlisting}
sudo systemctl start docker
\end{lstlisting}

\subsection{Compose not found}
Make sure you installed the Compose plugin:
\begin{lstlisting}
sudo apt install -y docker-compose-plugin
\end{lstlisting}

\subsection{Next step}
In the next chapter we will deploy PhotoPrism using Docker Compose, reusing the same principles:
\textbf{clear volumes}, \textbf{documented configuration}, and \textbf{easy upgrades}.